\scanpage{211}
\mspage{3}
% archive sheet number: 210

$X_{P_i}\simeq[1,i]_{P_i}$. Les $P_i\longrightarrow e$ forment une famille \deleted{épimorphique} de descente effective universelle \deleted{pour} dans $\mathscr T$, donc aussi dans $\mathscr T_{\unsure{\mathrm{revf}|\mathrm{rev}\,f|\mathrm{cf}}}$ (N.B. $\mathscr T_{\mathrm{revf}}$ \deleted{\unsure{admet des|stable}}\inserted{est stable} par produits finis -- en fait par $\underleftarrow{\lim}$ et $\underrightarrow{\lim}$ \uline{finis}).

\uline{Démonstration de B}

\deleted{En fait, on constate aisément que} \deleted{les conditions de A(i) impliquent G0 G1 G2 G3.}

\deleted{L'équivalence de (i) \unsure{:|à|et} (iii) tient au fait que} \deleted{les décompositions en sommes de $e$, dans $\mathscr T_{\mathrm{revf}}$} \deleted{ou dans $\mathscr T$, sont les \unsure{m\textsuperscript{es}} ; donc dire que $e$} \deleted{est somme finie d'objets connexes dans $\mathscr T$, ou} \deleted{dans $C$, revient au \unsure{m\textsuperscript{e}}.}

\deleted{Pour prouver que ces conditions impliquent (iv)} \deleted{il \deleted{faut}\inserted{reste à} prouver \deleted{que} qu'elles impliquent G4,} \deleted{i.e. tout revêtement fini d'un objet connexe $e$} \deleted{est somme finie d'objet connexes. En fait,} \deleted{si $X$ est de rang $n$ sur $e$, on trouve que} \deleted{le nb de composantes connexes de $X$ est $\le n$,} \deleted{et égal à $n$ ssi $X$ est trivial : c'est immédiat} \deleted{par récurrence sur $n$, et additivité du rang.}

\deleted{$(\mathrm{v})\Longrightarrow(\mathrm{iii})$ est immédiat ; on prend pour $\mathscr T$ le topos} \deleted{``somme'' des $B_{\pi_i}$. ($B_{\pi_i}=$ ensembles discrets sur} \deleted{lequel $\pi$ opère continûment).}

Il resterait à prouver $(\mathrm{ii})\Longrightarrow(\mathrm{iii})$. Je préfère prouver d'abord $(\mathrm{ii})\Longrightarrow(\unsure{\mathrm{i}|\mathrm{ii}|\mathrm{iii}}\deleted{\unsure{v}})$.

